% ==================== 附录 A ====================
\section{Origin 计算模板公式清单}

\subsection{列布局}

表 \ref{tab:calc-columns} 给出“晶体学计算器”工作表的推荐列布局。列名直接作为公式中的变量名，因此应避免使用中文、空格和重复名称。表 \ref{tab:diff-columns} 给出衍射图工作表的推荐列布局；如果实际表格的列顺序不同，只需把公式中的列字母换成对应的列即可。

\begin{table}[H]
  \centering
  \caption{晶体学计算器工作表的列布局}
  \label{tab:calc-columns}
  \small
  \setlength{\tabcolsep}{4pt}
  \begin{tabular}{cccccccc}
    \toprule
    列 & 名称 & 含义 & 列 & 名称 & 含义 & 列 & 名称/含义 \\
    \midrule
    A & a & 晶格常数 $a$ & H & a* & $a^*$ & O & k1 \\
    B & b & 晶格常数 $b$ & I & b* & $b^*$ & P & l1 \\
    C & c & 晶格常数 $c$ & J & c* & $c^*$ & Q & h2 \\
    D & alpha & $\alpha$ & K & alpha* & $\alpha^*$ & R & k2 \\
    E & beta & $\beta$ & L & beta* & $\beta^*$ & S & l2 \\
    F & gamma & $\gamma$ & M & gamma* & $\gamma^*$ & T & r1 \\
    G & V & 晶胞体积 & N & h1 & $h_1$ & U & r2 \\
    -- & -- & -- & -- & -- & -- & V & theta \\
    \bottomrule
  \end{tabular}
\end{table}

\begin{table}[H]
  \centering
  \caption{衍射图工作表的推荐列布局}
  \label{tab:diff-columns}
  \small
  \setlength{\tabcolsep}{4pt}
  \begin{tabular}{cccccccc}
    \toprule
    列 & 名称 & 含义 & 列 & 名称 & 含义 & 列 & 名称/含义 \\
    \midrule
    A & r1 & $r_1$ & G & y & 纵坐标 & M & k1 \\
    B & r2 & $r_2$ & H & I & 强度 & N & l1 \\
    C & theta & $\theta$ & I & h & $h$ & O & h2 \\
    D & m & $m$ & J & k & $k$ & P & k2 \\
    E & n & $n$ & K & l & $l$ & Q & l2 \\
    F & x & 横坐标 & L & h1 & $h_1$ & -- & -- \\
    \bottomrule
  \end{tabular}
\end{table}

\subsection{晶胞体积与倒易参数}

选中 G 列，按 \texttt{Ctrl+Q} 输入：
\begin{lstlisting}
A*B*C*sqrt(1-cos(D)^2-cos(E)^2-cos(F)^2+2*cos(D)*cos(E)*cos(F))
\end{lstlisting}
选中 H--M 列，分别输入：
\begin{lstlisting}
B*C*sin(D)/G
A*C*sin(E)/G
A*B*sin(F)/G
acos((cos(E)*cos(F)-cos(D))/(sin(E)*sin(F)))
acos((cos(D)*cos(F)-cos(E))/(sin(D)*sin(F)))
acos((cos(E)*cos(D)-cos(F))/(sin(E)*sin(D)))
\end{lstlisting}
注意：D、E、F 为角度，K、L、M 的输出也是角度；务必把 Origin 的 Angular Unit 设为 Degree。

\subsection{倒易矢量长度与夹角}

选中 T、U、V 列，分别输入：
\begin{lstlisting}
SQRT((N*H)^2+(O*I)^2+(P*J)^2+2*N*O*H*I*cos(M)+2*N*P*H*J*cos(L)+2*O*P*I*J*cos(K))
SQRT((Q*H)^2+(R*I)^2+(S*J)^2+2*Q*R*H*I*cos(M)+2*Q*S*H*J*cos(L)+2*R*S*I*J*cos(K))
ACOS((N*Q*H^2+O*R*I^2+P*S*J^2+(N*R+O*Q)*H*I*cos(M)+(N*S+P*Q)*H*J*cos(L)+(O*S+P*R)*I*J*cos(K))/(T*U))
\end{lstlisting}
式中的 N--S 为两组晶面指数，H--M 为倒易参数，T、U 为 $r_1,r_2$。如果列布局与表 \ref{tab:calc-columns} 不同，应逐项替换列字母，不要直接复制。

\subsection{衍射点坐标与指数}

在衍射图工作表中，按表 \ref{tab:diff-columns} 的布局输入：
\begin{lstlisting}
x = D*A + E*B*COS(C)
y = E*B*SIN(C)
h = D*L + E*O
k = D*M + E*P
l = D*N + E*Q
\end{lstlisting}
其中 A、B、C 为 $r_1,r_2,\theta$，D、E 为 $m,n$，L--Q 为两组基矢指数。$x,y$ 的单位是 \AA$^{-1}$；若要与实验花样的 $R$ 值比较，还需乘以相机长度 $L$ 或相机常数。

\subsection{结构因子}

在结构因子工作表中，按 A=$f$、B=$x$、C=$y$、D=$z$、E=$h$、F=$k$、G=$l$、H=phase、I=real、J=imag、K=intensity 的布局输入：
\begin{lstlisting}
phase = 360*(E*B + F*C + G*D)
real  = A*COS(H)
imag  = A*SIN(H)
I     = (TOTAL(I))^2 + (TOTAL(J))^2
\end{lstlisting}
每个衍射点占 8 行（对应表 \ref{tab:si-atoms} 的 8 个原子）。\texttt{TOTAL} 对这 8 行求和；若一次计算多个点，应把每个点的 8 行作为一个块，用 \texttt{SUM} 指定范围，或增加斑点编号列后用 \texttt{TOTALIF} 分组。强度为 0 或接近 $10^{-29}$ 的点即为消光点。

\subsection{自检值}

表 \ref{tab:self-check} 给出了用单晶硅（$a=5.43$~\AA）检查模板的期望值。若计算结果与表中数值不符，应先检查角度单位、列引用和公式中的交叉项。

\begin{table}[H]
  \centering
  \caption{单晶硅计算模板的自检值（$a=5.43$~\AA，Cu K$\alpha_1$）}
  \label{tab:self-check}
  \small
  \begin{tabular}{cccc}
    \toprule
    量 & 期望值 & 量 & 期望值 \\
    \midrule
    $V$ & 160.103~\AA$^3$ & $a^*$ & 0.18416~\AA$^{-1}$ \\
    $d_{111}$ & 3.1350~\AA & $r_{111}$ & 0.3189~\AA$^{-1}$ \\
    $d_{220}$ & 1.9198~\AA & $r_{220}$ & 0.5209~\AA$^{-1}$ \\
    $d_{400}$ & 1.3575~\AA & $r_{400}$ & 0.7366~\AA$^{-1}$ \\
    $2\theta_{111}$ & 28.45$^\circ$ & $2\theta_{220}$ & 47.31$^\circ$ \\
    \bottomrule
  \end{tabular}
\end{table}

\subsection{输入检查清单}

\begin{enumerate}[label=（\arabic*）,leftmargin=2.2em,itemsep=1pt]
  \item 角度单位是 Degree，不是 Radian；
  \item A--F 列的 $a,b,c,\alpha,\beta,\gamma$ 与目标物相一致；
  \item H--M 列的公式引用了正确的列，且分母为 $V=G$；
  \item T--V 列的分母为 $r_1r_2$，交叉项没有漏项；
  \item 衍射图表格中 $m,n$ 的范围足够大，能铺满绘图区；
  \item 每个衍射点的结构因子使用完整的 8 行原子坐标；
  \item \texttt{TOTAL} 或 \texttt{SUM} 的范围没有跨到相邻衍射点；
  \item $2\theta$ 是由 $d$ 通过 Bragg 方程得到的，不是把 $\theta$ 直接当成 $2\theta$；
  \item 用表 \ref{tab:self-check} 的 $V$、$a^*$、$d_{111}$、$2\theta_{111}$ 做一次自检。
\end{enumerate}

\section{图片素材索引}

本文正文使用的 56 张操作截图均来自 \texttt{Microstructure/新建文件夹}，按图号索引于表 \ref{tab:image-index}。每张截图在正文中只出现一次，未作裁剪或内容修改；为使文件路径兼容 XeLaTeX，复制到 \texttt{assets/} 时改用了 ASCII 文件名。其中图 4-5～4-9 的 (c) 二维散点图取自各带轴文件夹中的 \texttt{image.png}（Origin 重新导出的完整版），其余各图仍取自同名中文命名的截图。

\begin{table}[H]
  \centering
  \caption{正文图号与截图素材的对应关系}
  \label{tab:image-index}
  \small
  \setlength{\tabcolsep}{6pt}
  \renewcommand{\arraystretch}{1.02}
  \begin{tabular}{ccl}
    \toprule
    图号 & 张数 & 素材编号/内容 \\
    \midrule
    图 2-1 & 3 & c01--c03：新建工作簿、增加列、建立列名 \\
    图 2-2 & 2 & c04--c05：修正列名、输入公式 \\
    图 2-3 & 3 & c06--c08：角度单位、晶格参数、体积公式 \\
    图 2-4 & 2 & c09--c10：体积结果、倒易参数 \\
    图 2-5 & 2 & c11--c12：输入晶面、$r_1/r_2/\theta$ 结果 \\
    图 3-1 & 5 & p01--p05：VESTA 导入、分数坐标、导出 P1、记事本坐标 \\
    图 3-2 & 4 & p16--p19：结构因子、原子坐标、相位、强度 \\
    图 3-3 & 2 & p20--p21：110 输入、强度消光 \\
    图 4-1 & 3 & p06--p08：基矢计算、$r_1/r_2/\theta$、二维点阵 \\
    图 4-2 & 4 & p09--p12：新建表、列名、$m$ 和 $n$ 的范围 \\
    图 4-3 & 3 & p13--p15：$x/y$ 公式、$I=100$、当前 $hkl$ \\
    图 4-4 & 3 & p22--p24：Waterfall 设置、三维示例、二维示例 \\
    图 4-5 & 4 & zone001 的 (a)--(d) \\
    图 4-6 & 4 & zone111 的 (a)--(d) \\
    图 4-7 & 4 & zone332 的 (a)--(d) \\
    图 4-8 & 4 & zone334 的 (a)--(d) \\
    图 4-9 & 4 & zone345 的 (a)--(d) \\
    \midrule
    合计 & 56 & -- \\
    \bottomrule
  \end{tabular}
\end{table}

    \section{PDF 卡片绘制程序（Python）}
    \label{app:pdfcard-code}

    \subsection{构建思路}

    图 \ref{fig:pdf-card} 与表 \ref{tab:pdf-card} 由 \texttt{code/pdfcard.py} 生成。程序只依赖 Python 3 标准库和 \texttt{matplotlib}，整个流程分为六步，与代码中的函数一一对应：

    \begin{enumerate}[label=（\arabic*）,leftmargin=2.2em,itemsep=2pt]
      \item \textbf{固定输入}：晶格常数 $a=5.43$~\AA、Cu K$\alpha_1$ 波长 $\lambda=1.5406$~\AA、常数原子散射因子 $f_0=14$、温度因子 $B=0.5$~\AA$^2$，以及 Fd-3m 的 8 个 Si 原子分数坐标。
      \item \textbf{枚举独立反射}：只遍历 $0\le l\le k\le h\le 8$。任意交换或改变正负号得到的 $hkl$ 与它同属一个晶面族，$d$ 和 $|F|$ 完全相同，其贡献通过多重性统一计入，因此不必重复计算。
      \item \textbf{结构因子与消光}：按 $F/f=\sum_j\exp[2\pi i(hx_j+ky_j+lz_j)]$ 求和后取模。取模后小于 $10^{-6}$ 的晶面即 $|F|^2=0$，判为消光并直接跳过；同时跳过 $\lambda/(2d)>1$、超出 Cu K$\alpha_1$ 衍射极限的反射。
      \item \textbf{几何量与多重性}：立方晶系用 $d=a/\sqrt{h^2+k^2+l^2}$，由 $2\theta=2\arcsin[\lambda/(2d)]$ 得衍射角；多重性 $N$ 取“排列数 $\times2^{\text{非零指数个数}}$”，等价于把交换和正负号都算进去。Lorentz-偏振因子按 $LP=(1+\cos^22\theta)/(\sin^2\theta\cos\theta)$ 计算。
      \item \textbf{两种强度模型与分组}：$I\propto N|F|^2LP$。模型一取 $f=f_0=14$；模型二把 $f$ 换成 Cromer--Mann 的 $f(s)$ 并乘温度因子 $e^{-2Bs^2}$。同一个 $d$ 值下可能对应多组 $hkl$（如 333 与 511），用 $\mathrm{round}(d,5)$ 作为键把它们合并，组内强度与多重性分别相加。
      \item \textbf{归一化与绘图}：以每套模型的最大强度为 100\%，用 \texttt{vlines} 在 $2\theta$ 位置画棒，双面板共享纵轴，300~dpi 保存为 \texttt{assets/pdfcard.png}；同时把每行数据按“表 2-1 的列格式”打印到终端。
    \end{enumerate}

    这样组织的好处是：把“数据计算”和“绘图”解耦，\texttt{collect\_reflections()} 的输出既可以是表格行，也可以是画图用的数组；只要改动顶部的 $a$、$\lambda$、$f_0$、$B$ 或原子坐标，就能直接用于其它立方晶体。程序中的每一步都对应正文的式 \eqref{eq:structure-factor-simple}、式 \eqref{eq:bragg}、式 \eqref{eq:intensity-powder} 和式 \eqref{eq:cromer-mann}，因此可以逐项核对。

    \subsection{完整代码}

    \lstinputlisting[language=Python,basicstyle=\ttfamily\scriptsize,
      caption={单晶硅 PDF 卡片的计算与绘图脚本 \texttt{code/pdfcard.py}},
      label={lst:pdfcard}]{code/pdfcard.py}

    \subsection{运行与自检}

    在 \texttt{课程作业} 目录下执行
    \begin{lstlisting}
    cd code
    python pdfcard.py
    \end{lstlisting}
    程序会在终端打印 12 行可直接粘贴到表 \ref{tab:pdf-card} 的数据，并在 \texttt{../assets/pdfcard.png} 生成图 \ref{fig:pdf-card}。若缺少绘图库，先执行 \texttt{pip install matplotlib}。自检要点：
    \begin{enumerate}[label=（\arabic*）,leftmargin=2.2em,itemsep=1pt]
      \item $f(0)=6.2915+3.0353+1.9891+1.5410+1.1407\approx14$，与 Si 的原子序数一致；
      \item 111 反射 $d=3.1350$~\AA、$2\theta=28.45^\circ$、$|F|^2=32f^2=6272$，相对强度为 100；
      \item 110、200、222、420、600 等不出现（被消光判据剔除），与第 3 章表 \ref{tab:extinction-examples} 一致；
      \item 脚本打印的 12 行数值应与 Origin 模板逐行相同，若不同，优先检查原子坐标的小数位和 $\lambda$ 的取值。
    \end{enumerate}

